% !TEX root = ../main.tex
\section{Discussion}

Across the meta-plots the two-stage estimator reproduced the truncated fits and most often delivered the lowest AICc among the three alternatives. The complete single-stage estimator remained biased near the truncation limits, whereas the truncated single-stage (1st) and two-stage complete (2sc) solutions were visually similar in Figure~\ref{fig:comparison}. These findings show that the two-stage workflow typically matches the reference truncated procedure without bespoke truncated-density implementations, while acknowledging a handful of cases where 1sc is marginally preferred by AICc.
We emphasize that 1sc is included as a deliberate operational benchmark rather than a strawman comparator: it reflects the common complete-form fitting pattern applied to truncated grouped stand tables when truncation is not modelled explicitly.

Two meta-plots (white pine--mixedwood and red pine--mixedwood) produce unusually large scale parameters in Table~\ref{tab:comparison}. Their truncated histograms are nearly flat, so the optimiser pushes shape parameters below one and inflates the scale to hold the area to unity. Both 1st and 2sc converge on the same regime, indicating a data-driven effect rather than a two-stage weakness.

We limited the candidate set to the two-parameter Weibull and Gamma distributions because they cover common stand structures with few parameters, both arise as special cases of the generalised gamma family, and a compact family keeps attention on the methodological swap between the truncated and two-stage estimators. A three-parameter Weibull (with location) is an alternative way to handle left truncation; we treat it as future work because it changes the parameterisation and increases flexibility in ways that would blur the direct 1st-vs-2sc comparison in a brief format.

A companion repository (URL withheld for double-blind review) documents the tallies, weighting choices, and scripts that regenerate the analysis. Robustness summary tables are provided in the manuscript appendix (Tables~\ref{tab:app_shape_wins}--\ref{tab:app_sim_summary}). With that scaffold, 2sc matches or improves on the truncated baseline and outperforms the unadjusted 1sc baseline in most tested settings. Because the interpretation of $b$ is application-specific, we view the two-stage workflow as an alternative that should be stress-tested under other truncation rules or distributional shapes before being treated as a universal replacement. The current analysis focuses on inverse-J PSP stand tables; other shapes warrant future study.

The robustness analysis includes empirical shape-stratified summaries across all 32 PSP combinations and a canonical synthetic benchmark spanning inverse-J, unimodal-mid, near-flat, and right-shifted cases. This evidence improves coverage of plausible stand-structure conditions, but it also confirms that method preference is context-dependent in low-information regimes. Accordingly, we frame 2sc as a strong practical option under tested conditions rather than a universal replacement.
